% FILE: 04_implementation_surface.tex
% Project: lifecycle
% Role: process-lifecycle-state-definitions-and-transitions

\section{Implementation Surface}
\label{sec:lifecycle-implementation}

\subsection{Source Files}
\label{sec:lifecycle-implementation-sources}

The lifecycle implementation surface consists of 43 Markdown files authored between
31 May and 1 June 2026. The files are organized into four functional groups.

The six operational stage definition files are the core of the implementation:
\texttt{define\_design-state\_process\_intelligence.md},
\texttt{define\_simulation-state\_process\_intelligence.md},
\texttt{define\_monitoring-state\_process\_intelligence.md},
\texttt{define\_repair-state\_process\_intelligence.md},
\texttt{define\_optimization-state\_process\_intelligence.md}, and
\texttt{define\_decommission-state\_process\_intelligence.md}. Each file specifies the
stage's MAPE-K element, input event or log type, analysis engine or algorithm, planning
or optimization method, execution controller, transition class, and knowledge base asset.
Companion narrative files (\texttt{design.md}, \texttt{simulation.md},
\texttt{monitoring.md}, \texttt{repair.md}, \texttt{decommissioning.md}) provide
extended prose exposition of each stage's formal content.

The seven supporting state files cover the non-operational lifecycle phases:
\texttt{ACQUISITION.md} (M\&A due diligence ingestion and the diligence gap),
\texttt{ARCHIVE.md} (cryptographic archival and cold data strategies),
\texttt{BOARD\_PROJECTION.md} (executive dashboard translation rules), and
definition files for Construction, Activation, Integration, and Operation states.

The autonomic control files are \texttt{MAPE\_K\_MAP.md},
\texttt{define\_blue\_river\_dam\_lifecycle\_gate\_map.md},
\texttt{define\_autonomic\_knowledge\_actuation\_map.md},
\texttt{define\_final\_receipt\_state.md}, and
\texttt{board-projection-holography.md}. These govern the feedback orchestration
protocols and the quality gate criteria.

The five taxonomy files---\texttt{define\_process\_asset\_taxonomy.md},
\texttt{define\_process\_debt\_taxonomy.md},
\texttt{define\_process\_readiness\_taxonomy.md},
\texttt{define\_process\_residual\_taxonomy.md}, and
\texttt{define\_process\_risk\_taxonomy.md}---plus the
\texttt{define\_false-claim\_taxonomy.md} and
\texttt{define\_false-claim\_taxonomy.md} files complete the classification surfaces.
The maturity model (\texttt{MATURITY\_MODEL.md}) and the completeness audit
(\texttt{audit\_\_lifecycle\_completeness.md}) close the file inventory.

\subsection{Commands}
\label{sec:lifecycle-implementation-commands}

The lifecycle project does not define a build tool or CLI command surface of its own.
It is a doctrine and specification layer, not an execution layer. The execution commands
that operationalize the lifecycle---token replay, alignment computation, Inductive Miner
discovery, receipt emission---are obligations assigned to the \texttt{wasm4pm} execution
engine, which is documented as a separate downstream product.

The Python assertion blueprints in \texttt{checkpoint\_\_lifecycle\_model\_complete.md}
constitute the specification of what commands must be executable by any conforming
implementation:
\begin{itemize}
  \item \texttt{assert\_soundness(net)}: verify WF-net structure, 1-boundedness, liveness,
    and absence of dead transitions.
  \item \texttt{assert\_fitness(net, log, threshold=0.95)}: compute alignment fitness and
    assert it meets the threshold.
  \item \texttt{assert\_no\_ghost\_transitions(net, log)}: verify every labeled transition
    is fired at least once in the log or is a silent transition.
  \item \texttt{assert\_decommission\_receipt(receipt, pub\_key)}: verify the
    cryptographic decommissioning receipt structure and signature.
\end{itemize}

These blueprints are specification artifacts, not executable test files. No \texttt{.py},
\texttt{.rs}, or test runner files exist within the \texttt{lifecycle/} directory. The
gap between specification and execution is documented explicitly in the checkpoint status.

\subsection{Data and Ontology Surfaces}
\label{sec:lifecycle-implementation-data}

The lifecycle project does not include TTL, RDF, or OWL ontology files. The type-level
ontology of the lifecycle is expressed through the \texttt{wasm4pm-compat} Rust type
system, with const generics encoding soundness classes and typestate methods enforcing
object lifecycle phase transitions. The Object Lifecycle typestates are formally defined
as $\mathrm{PHASE} \in \{\texttt{Created}, \texttt{Active}, \texttt{Modified},
\texttt{Archived}, \texttt{Deleted}\}$ with intermediate type aliases introduced to
prevent nightly Rust compiler E0391 variance cycle errors.

The event log data formats recognized by the lifecycle are XES (IEEE~1849) and OCEL~2.0.
OCEL~2.0 is preferred for multi-object processes because it preserves exact
event-to-object relationships that would be distorted by flattening to a single-case
XES structure. The ACQUISITION.md document specifies that admitted OCEL logs covering
a 24-month due diligence period constitute the Layer 2 evidence requirement for
M\&A process intelligence packages.

Receipt data is stored in JSON format in the parent \texttt{receipts/} directory. The
canonical receipt fields are: model hash (\texttt{BLAKE3(N)}), log hash
(\texttt{BLAKE3(L\_final)}), total cost ($C_{\mathrm{total}}$), final fitness
($F_{\mathrm{final}}$), and retirement timestamp ($T_{\mathrm{retire}}$), signed with
Ed25519 using the governance private key.

\subsection{Templates and Rendered Artifacts}
\label{sec:lifecycle-implementation-templates}

The Autonomic Knowledge Actuation Map (\texttt{define\_autonomic\_knowledge\_actuation\_map.md})
provides a canonical tabular template for specifying each stage's complete actuation
profile. The template columns are: Stage, MAPE-K Element, Input Event or Log, Analysis
Engine or Algorithm, Planning or Optimization method, Execution Controller, Transition
Class, and Knowledge Base Asset. This template is the reference form that any conforming
lifecycle stage specification must populate.

The MAPE-K mapping matrix in \texttt{MAPE\_K\_MAP.md} provides a canonical tabular
template for the \texttt{wasm4pm} capability requirements per MAPE-K function, with
columns for MAPE-K function, \texttt{wasm4pm} capability, and \texttt{wasm4pm-compat}
type. The matrix explicitly marks each capability as either present in
\texttt{wasm4pm-compat} (the structure layer) or required in \texttt{wasm4pm} (the
execution layer), making the implementation gap visible as a structured artifact rather
than a prose disclaimer.

The Process Intelligence Diligence Package template, defined in \texttt{ACQUISITION.md},
specifies five layers of evidence that a seller operating at L3 or above can produce
for M\&A due diligence: model evidence (Layer 1), log evidence (Layer 2), conformance
evidence (Layer 3), repair evidence (Layer 4), and replay verification (Layer 5). This
template is the primary manufactured artifact of the lifecycle framework for M\&A
transaction contexts.
